\documentclass[11pt]{article}
% Supplementary material for "Discovering and Certifying Partial Symbolic
% Invariants for Neuro-Symbolic Multi-Agent World Models" (AAAI-28 working
% draft). This document is NOT page-limited and is intended to accompany
% paper.tex as the Code and Data / Supplementary Document at submission time.

\usepackage[margin=1in]{geometry}
\usepackage{amsmath, amssymb, amsfonts}
\usepackage{amsthm}
\usepackage{bm}
\usepackage{booktabs}
\usepackage{array}
\usepackage{multirow}
\usepackage{mathtools}
\usepackage{xcolor}
\usepackage{algorithm}
\usepackage{algorithmic}
\usepackage{hyperref}
\usepackage[capitalize,noabbrev]{cleveref}
\usepackage{natbib}

\newcommand{\E}{\mathbb{E}}
\newcommand{\R}{\mathbb{R}}
\newcommand{\ph}{\textcolor{red}{\texttt{[TBD]}}}

\theoremstyle{plain}
\newtheorem{theorem}{Theorem}[section]
\newtheorem{proposition}[theorem]{Proposition}
\newtheorem{lemma}[theorem]{Lemma}
\newtheorem{corollary}[theorem]{Corollary}
\theoremstyle{definition}
\newtheorem{definition}[theorem]{Definition}
\newtheorem{assumption}[theorem]{Assumption}
\theoremstyle{remark}
\newtheorem{remark}[theorem]{Remark}

\title{Supplementary Material\\[2pt]\large Discovering and Certifying Partial Symbolic Invariants for Neuro-Symbolic Multi-Agent World Models}
\author{Anonymous Submission}
\date{}

\begin{document}
\maketitle

\tableofcontents

\section{Scope of this document}

This supplementary material accompanies the AAAI-28 working draft \texttt{paper.tex}. It contains: the full transport catalog and estimators (\autoref{app:dsl}); the complete discovery algorithm, its data-split discipline, and complexity (\autoref{app:algo}); full proofs of Theorems 5.1--5.5 and the composition remark, including the constructive Dec-POMDP counterexample for the non-identifiability theorem (\autoref{app:proofs}); the extended experimental protocol, baseline design space, and ablation list (\autoref{app:protocol}); the LLM-proposer prompting protocol (\autoref{app:llm}); and reproducibility notes (\autoref{app:repro}). Numbers marked \ph{} are structural placeholders, not measurements.

\section{Transport catalog and precondition language}
\label{app:dsl}

\subsection{Data partitions}
Discovery uses four disjoint partitions, never mixed across roles:
\begin{itemize}
  \item $\mathcal{D}_{\mathrm{prop}}$: structure search (which $(\varphi,c,\rho)$ to propose). Never used to certify.
  \item $\mathcal{D}_{\mathrm{fit}}$: parameter/threshold estimation ($\delta_a$, $[\ell,u]$, decision thresholds $\theta$). Never used to certify.
  \item $\mathcal{D}_{\mathrm{test}}$: falsification only (F1). Never used to select or tune candidates.
  \item $\mathcal{D}_{\mathrm{shift}}$: held-out environment variants and alternative data generators (F3, F4).
\end{itemize}
This separation makes the union bound in \autoref{thm:t2-full} apply to the number of candidates \emph{proposed}, $\hat M = |\hat{\mathcal{R}}|$, rather than to the size of the catalog's search space, and it is what allows $\mathcal{D}_{\mathrm{test}}$ to be treated as i.i.d.\ for \autoref{thm:t1-full}: we recommend sub-sampling one transition per episode when transitions within an episode are not i.i.d.\ (the cheap, conservative option); a block/martingale argument is a looser but more sample-efficient alternative left as a remark, not adopted in the main protocol.

\subsection{Readings $\varphi$}
Three levels, chosen because they are closed under permutation of feature indices --- the property that makes the anonymization experiment (\autoref{app:protocol}) a test of the algorithm rather than of feature-name leakage:
\begin{itemize}
  \item \textbf{Atomic}: $f_{i,k}\mapsto \varphi\in\R$, $|\Phi_1|=d=n\cdot n_f$.
  \item \textbf{Relational}: $f_{i,k}-f_{j,l}\mapsto\varphi$, $|\Phi_2|=O(d^2)$.
  \item \textbf{Aggregated}: $\sum_i f_{i,k}$, $\max_i f_{i,k}$, $\#\{i:f_{i,k}=v\}$, restricted to the agent axis only, $O(n_f)$. This restriction keeps $|\Phi|$ dominated by $d^2$ and is exactly the form multi-agent conservation laws take (total resource, count of live units).
\end{itemize}

\subsection{Transports $\rho$}
\begin{itemize}
  \item \textbf{Family A --- bounded affine}: $\rho(\varphi,a)=\mathrm{clip}(\alpha\varphi+\delta_a,\ell,u)$, $\alpha\in\{0,1\}$. Closed-form fit on $\mathcal{D}_{\mathrm{fit}}$: $\delta_a=\mathrm{mode}(\varphi_{t+1}-\varphi_t\mid a)$, $[\ell,u]=[\min,\max]$ observed; no gradient required. $(\alpha,\delta_a,\mathrm{clip})=(1,0,\text{no})$ recovers persistence; $(0,v,\text{no})$ a conditional constant; $(1,\neq0,\text{no})$ free motion/increment; $(1,\neq0,\text{yes})$ cooldowns, health bounds, grid boundaries.
  \item \textbf{Family B --- delayed copy}: $\rho=\psi_{t-m}$ for another reading $\psi$ and $m\in\{0,\dots,m_{\max}\}$. Covers pickup/drop and transfer effects; type compatibility (cardinality, integrality) is inferred from value domains, never from feature names.
  \item \textbf{Family C --- slot permutation}: $\varphi^{\sigma(i)}_{t+1}=\varphi^i_t$. Covers agent-slot reassignment; optional in a first implementation pass since it is expensive and rarely decisive relative to Families A/B.
\end{itemize}

\subsection{Preconditions $c$}
Conjunctions of at most $\beta$ atoms (we recommend $\beta=2$; $\beta=3$ increases $\hat M$ substantially for a gain that has not been demonstrated and should be tested as an ablation), plus $c=\top$. Admissible atoms: action literal ($a^i=u$), threshold test ($\varphi'\Join\theta$, $\theta$ taken from observed deciles), categorical equality, and relational atom ($\varphi_i'=\varphi_j'$).

\subsection{Non-vacuity screen}
\label{app:screen}
Applied on $\mathcal{D}_{\mathrm{prop}}$ before falsification, this is the ``coverage'' side of a precision/coverage trade-off and the easiest step to omit incorrectly:
\begin{itemize}
  \item reject if $\varphi$ is near-constant;
  \item reject if the rate of change is $\approx 0$, i.e., $\Pr[\varphi_{t+1}\neq\varphi_t]\approx0$ under the identity transport --- a persistence rule on a feature that never changes is vacuously true and inflates $\hat\kappa$ without adding information;
  \item reject if $\Pr[c] < \pi_{\min}$ (insufficient support);
  \item reject if removing the candidate as a covered feature does not measurably degrade residual prediction accuracy on $\mathcal{D}_{\mathrm{fit}}$ (the correlation criterion referenced in the main paper's injection section).
\end{itemize}
Retained candidates lie on the Pareto frontier of $(1-\hat\varepsilon,\ \mathrm{cov})$ where $\mathrm{cov}(r)=\Pr[c]\cdot H(\varphi(\omega_{t+1})\mid c)/H_{\max}$.

\section{Discovery algorithm: full specification}
\label{app:algo}

\begin{algorithm}[H]
\caption{Discover, falsify, inject (full)}
\begin{algorithmic}[1]
\STATE \textbf{Stage 0 (screening).} For each feature $k$, fit a small predictor of $f_{t+1,k}$ on $\mathcal{D}_{\mathrm{prop}}$ and rank by residual error; this only ranks candidacy for coverage, it does not itself produce a rule.
\STATE \textbf{Stage 1 (propose).} Compare three proposers in ablation:
  \begin{itemize}
    \item \textbf{P1}: enumerate the catalog (\autoref{app:dsl}) up to precondition depth $\beta$. Baseline; gives $\hat M$ in closed form for \autoref{thm:t2-full}.
    \item \textbf{P2}: fit a decision tree on $f_{t+1,k}$ over $\mathcal{D}_{\mathrm{prop}}$; extract pure leaves as (precondition, effect) pairs. $\hat M=$ number of retained leaves.
    \item \textbf{P3}: prompt an LLM with anonymized features and example transitions (\autoref{app:llm}); $\hat M=$ number of syntactically valid candidates returned.
  \end{itemize}
\STATE \textbf{Stage 2 (falsify)}, four gates of increasing severity, each eliminatory:
  \begin{itemize}
    \item \textbf{F1 (passive)}: holdout transitions from $\mathcal{D}_{\mathrm{test}}$.
    \item \textbf{F2 (active)}: exploration directed at states satisfying $c$ --- direct environment reset in GridCraft; rejection sampling from a diverse behavioral-policy pool in Overcooked/SMACv2. This is what distinguishes the protocol from passive validation (and from WALL-E-style induce/prune).
    \item \textbf{F3 (cross-variant)}: environment variants, testing the invariance criterion directly.
    \item \textbf{F4 (cross-generator)}: data from random, scripted, and alternative MARL-trained generators; a direct empirical attack on the policy-shift degradation result (Theorem~5.4 in the main paper; \autoref{thm:t4-full} here).
  \end{itemize}
  A rule is certifiable only if $N_c(r)\ge N_{\min}$; rules with rare preconditions are reported as \emph{non-certifiable}, not false.
\STATE \textbf{Memory-order test.} Each candidate is tested at $m\in\{0,1,2,\dots\}$; report the distribution of the minimal sufficient $m$ --- itself a result under Dec-POMDPs, and the empirical counterpart to the deferred-effect/latent-confounder distinction in \autoref{thm:t5-full}.
\STATE \textbf{Stage 3 (promote/retrograde).} Confidence $q_r=1-\hat\varepsilon_r$, $\hat\varepsilon_r=\ln(\hat M/\delta)/N_c(r)$ absent violations, Clopper--Pearson otherwise. $q_r<\tau_{\mathrm{soft}}\Rightarrow$ reject; $\tau_{\mathrm{soft}}\le q_r<\tau_{\mathrm{hard}}\Rightarrow$ soft, weighted $\lambda q_r$; $q_r\ge\tau_{\mathrm{hard}}\Rightarrow$ hard (projection or residual). Suggested starting values $\tau_{\mathrm{soft}}=0.9$, $\tau_{\mathrm{hard}}=0.99$, swept in ablation.
\STATE \textbf{Retrogradation.} An online monitor recomputes the bound with accumulated counters; on a counterexample, the rule is demoted and the feature returns to the neural target. Retrogradation is \emph{sticky}: a retrograded rule requires roughly double the evidence to re-promote, to avoid oscillation.
\STATE \textbf{Conflicts.} Two certified rules writing the same position: prefer the more specific precondition, then the higher $q$; if both are hard and contradictory, retrograde both and log the conflict as a diagnostic.
\end{algorithmic}
\end{algorithm}

\subsection{Complexity}
Let $d=nn_f$, $d_u=(1-\kappa)d$, $|\mathcal{R}|$ atomic rules, $C_\theta(d)$ the cost of one neural update with output dimension $d$, $H$ the rollout horizon. Rule evaluation is $O(|\mathcal{R}|q)$ time, $O(d)$ memory, for rules touching at most $q$ input features (bounded arity in our catalog, not $q=O(d)$). Per training transition: regularization and projection cost $O(C_\theta(d)+|\mathcal{R}|q+d)$; residual dynamics costs $O(C_\theta(d_u)+|\mathcal{R}|q+d)$, reducing the neural output dimension by a factor of $(1-\kappa)$ for architectures whose per-step cost is linear in output width. An $H$-step open-loop rollout multiplies these by $H$, with rollout-storage memory $O(Hd)$ if fully logged or $O(d+h)$ if streamed ($h$ = hidden-state size).

\section{Full proofs}
\label{app:proofs}

\subsection*{Notation}
Dec-POMDP $\mathcal{M}$, behavioral policy $\pi$ inducing $d^\pi$ over transitions $x=(h^j_{t-1},\omega^j_t,a^j_t,\omega^j_{t+1})$; $x^-=(h^j_{t-1},\omega^j_t,a^j_t)$. A rule $r=(\varphi,c,\rho)$ has violation indicator $V_r(x)=\mathbb{1}[c(x^-)\wedge \varphi(\omega_{t+1})\neq\rho_{a_t}(\varphi(\omega_t),h_{t-1})]$ and violation rate $\varepsilon^\pi_r=\Pr_{d^\pi}[V_r=1\mid c]$. $B=\mathrm{diam}(\mathcal{Y})$, $N_c(r)=|\{x\in\mathcal{D}_{\mathrm{test}}: c(x^-)\}|$. The fact making everything below licit: $c$ depends only on $x^-$, so conditioning on $N_c$ does not bias the violation indicators, and $\{x: c(x^-)\}$ is i.i.d.\ from the conditional law by construction of $\mathcal{D}_{\mathrm{test}}$.

\begin{assumption}[Certification data regularity, full statement]
\label{ass:regularity-full}
Fix $r$ independently of $\mathcal{D}_{\mathrm{test}}$. Let $(x_k)_{k\ge1}$ enumerate, in temporal order, the elements of $\mathcal{D}_{\mathrm{test}}$ satisfying $c$, and let $V_k=V_r(x_k)\in\{0,1\}$. Two regimes:
\begin{itemize}
  \item[(a)] \textbf{Sub-sampled.} At most one transition per episode is retained in $\mathcal{D}_{\mathrm{test}}$. Then $(x_k)$ are i.i.d.\ draws from $d^\pi(\cdot\mid c)$, and $(V_k)$ are i.i.d.\ Bernoulli($\varepsilon_r^\pi$); $N_c(r)$ is fixed by the (pre-registered) test set size.
  \item[(b)] \textbf{Any-time-valid.} Transitions within an episode may be correlated in arbitrary, unmodeled ways, and $N_c(r)$ may itself depend on the data (e.g., falsification stopped early, or pooled adaptively across gates). We only require that $(V_k)$ is a $\{0,1\}$-valued predictable sequence with $\Pr[V_k=1\mid\mathcal{F}_{k-1}]=\varepsilon_r^\pi$ for all $k$, where $\mathcal{F}_k=\sigma(x_1,\dots,x_k)$ is the natural filtration.
\end{itemize}
Regime (b) is the default; (a) is a cheaper practical fallback (implemented simply by discarding all but one transition per episode) when the supermartingale construction below is not implemented in the falsification pipeline.
\end{assumption}

\begin{theorem}[Certification by survival]
\label{thm:t1-full}
Under \autoref{ass:regularity-full}, if no violation is observed among $N_c(r)$ transitions satisfying $c$, then with probability at least $1-\delta$: $\varepsilon^\pi_r \le \big(\ln(1/\delta)+O(\log\log N_c(r))\big)/N_c(r)$ under regime (b), uniformly over $N_c(r)$ (including if $N_c(r)$ is a stopping time); under regime (a) the $O(\log\log N_c(r))$ term is exactly zero and the bound is $\ln(1/\delta)/N_c(r)$.
\end{theorem}
\begin{proof}
\textbf{Regime (a).} Conditional on $N_c=n$, the $n$ violation indicators are i.i.d.\ Bernoulli($\varepsilon_r$). The probability of observing zero violations is $(1-\varepsilon_r)^n\le e^{-\varepsilon_r n}$. If $\varepsilon_r>\ln(1/\delta)/n$, this probability is $<\delta$. Hence the event ``survive with $\varepsilon_r$ above the threshold'' has probability $<\delta$ conditional on $N_c$, hence unconditionally (since $n$ is fixed in advance under (a), there is no optional-stopping issue to control).

\textbf{Regime (b).} Fix a null value $\varepsilon$ and define, on the filtration $(\mathcal{F}_k)$ of \autoref{ass:regularity-full}, the sequential likelihood-ratio process
$$M_k(\varepsilon) = \prod_{j=1}^{k}\left(\frac{1-\varepsilon}{1-p_j}\right)^{\mathbb{1}[V_j=0]}\left(\frac{\varepsilon}{p_j}\right)^{\mathbb{1}[V_j=1]}, \qquad p_j := \Pr[V_j=1\mid\mathcal{F}_{j-1}].$$
Under the (null) hypothesis $\varepsilon_r^\pi\le\varepsilon$, $M_k(\varepsilon)$ is a nonnegative supermartingale with $M_0=1$ (a test supermartingale in the sense of Ville). By Ville's maximal inequality, $\Pr[\exists k: M_k(\varepsilon)\ge 1/\delta]\le\delta$. If zero violations are observed through $k=N_c(r)$, $M_{N_c(r)}(\varepsilon)=\prod_{j=1}^{N_c(r)}\frac{1-\varepsilon}{1-p_j}$; taking $\varepsilon=\ln(1/\delta)/N_c(r)$ recovers the fixed-sample exponent up to the boundary-crossing correction that any valid always-valid confidence sequence must pay for holding uniformly over $k$ (rather than at one pre-specified $k$), which is $O(\log\log N_c(r))$ for the mixture/Gamma-exponential constructions of \citet{howard2020timeuniform}; we do not re-derive the exact closed-form boundary here and refer to that paper's Theorem 1 and Table 3 for the explicit constant.
\end{proof}
For $\delta=0.05$ and regime (a): $\varepsilon_r\le 3/N_c$ (the ``rule of three''). With $k\ge1$ observed violations, replace by the Clopper--Pearson upper bound $\bar\varepsilon(k,N_c,\delta)$ under (a), or its any-time-valid analogue under (b)~\citep{howard2020timeuniform}.

\textbf{Why this matters for the rest of \autoref{app:proofs}.} Every downstream result (\autoref{thm:t2-full}, \autoref{cor:t2b-full}) is a union bound or expectation computed from \autoref{thm:t1-full}'s per-candidate tail probability; both inherit whichever regime of \autoref{ass:regularity-full} is in force without further modification, since the union bound and linearity-of-expectation arguments used below do not depend on which of the two proofs above supplied the per-candidate tail bound.

\subsection*{Effective budget}
Let $\hat M$ count proposed rules $(\varphi,c,\rho)$ before any grid over promotion thresholds or memory order. Since $\tau_{\mathrm{soft}},\tau_{\mathrm{hard}}\in\mathrm{grid}$, $m\in\{0,\dots,m_{\max}\}$, and precondition depth $\beta\in\mathrm{grid}$ are chosen (swept or fit) on the same $\mathcal{D}_{\mathrm{prop}}\cup\mathcal{D}_{\mathrm{fit}}$ used to propose candidates, treating them as free is equivalent to certifying only one member of a larger, implicitly-searched family. We instead certify the whole family by inflating the budget:
$$\hat M_{\mathrm{eff}} = \hat M\cdot|\mathrm{grid}(\tau_{\mathrm{soft}})|\cdot|\mathrm{grid}(\tau_{\mathrm{hard}})|\cdot(m_{\max}{+}1)\cdot|\mathrm{grid}(\beta)|.$$
This is deliberately a crude, conservative over-count (it does not exploit that many $(\tau,m,\beta)$ combinations yield the same certified rule set, or that gates are eliminatory so most of the grid is never actually evaluated for most candidates); a tighter accounting is possible but not necessary given the union bound's logarithmic dependence on the budget (below).

\begin{theorem}[Search cost, per gate]
\label{thm:t2-full}
Let $\mathcal{R}$ be the $\hat M_{\mathrm{eff}}$ candidates proposed from $\mathcal{D}_{\mathrm{prop}}\cup\mathcal{D}_{\mathrm{fit}}$ only, each tested at up to $K$ gates ($K=4$: F1--F4) on data disjoint across gates, with gate $k$ generated by policy $\pi_k$ ($\pi_1=\pi$). With probability at least $1-\delta$, simultaneously for every $r\in\mathcal{R}$ and every gate $k$ it reaches, $\varepsilon^{\pi_k}_r \le \ln(\hat M_{\mathrm{eff}} K/\delta)/N_c^{(k)}(r)$.
\end{theorem}
\begin{proof}
Apply \autoref{thm:t1-full} to each of the $\hat M_{\mathrm{eff}}\times K$ (candidate, gate) pairs at level $\delta/(\hat M_{\mathrm{eff}}K)$, then a union bound over all pairs. The union ranges over \emph{every} pair regardless of whether candidate $r$ is actually tested at gate $k$ in a given run of the discovery loop (Algorithm 1 in the main paper): since which pairs are realized is itself a (data-dependent, adaptive) function of the earlier gates' outcomes, taking the union over the full combinatorial space rather than only the realized subset is what makes the bound valid under this adaptivity --- a realized subset of a uniformly-controlled family inherits the same uniform guarantee, for any (even adversarially chosen) selection rule. Independence of each candidate from the data used to test it at any single gate is guaranteed by the partition discipline of \autoref{app:dsl} together with the fresh-data requirement on re-promotion below.
\end{proof}

\paragraph{Retrogradation and re-promotion.} A rule demoted by an online counterexample during training may be re-promoted, but only against \emph{fresh} falsification data disjoint from every previous test of that rule (including its original F1--F4 certification and any earlier re-promotion attempts), and at most $R_{\max}$ times. Both provisions are necessary: without the freshness requirement, repeated adaptive re-testing against previously-seen data is an uncontrolled multiple-comparisons problem outside the scope of \autoref{thm:t2-full}'s union bound; without the cap, $R_{\max}$ is unbounded and no finite budget correction covers it. With both, the effective budget for the full lifecycle of a single rule is $\hat M_{\mathrm{eff}}\times K\times R_{\max}$, and \autoref{thm:t2-full} applies verbatim with this product in place of $\hat M_{\mathrm{eff}}K$.

\autoref{thm:t2-full} and \autoref{cor:t2b-full} below are two readings of the \emph{same} tail bound, not competing arguments, and it is important not to conflate them (an earlier version of this derivation did). Read as an error bound at fixed sample size $N_{\min}$, the union-bound cost of a larger $\hat M_{\mathrm{eff}}$ is only $\ln\hat M_{\mathrm{eff}}$: going from $10^2$ to $10^6$ candidates only inflates the required sample budget by a factor $\ln(10^6)/\ln(10^2)\approx3$ (equivalently, $N_{\min}\ge\ln(\hat M_{\mathrm{eff}}/\delta)/\varepsilon^\star$ grows logarithmically in $\hat M_{\mathrm{eff}}$). Read instead as an expected \emph{count} of false discoveries at fixed $N_{\min}$, the count is linear in the number of bad candidates actually proposed:

\begin{corollary}[Expected false discoveries]
\label{cor:t2b-full}
Let $\mathcal{R}_{\mathrm{bad}}=\{r\in\mathcal{R}:\varepsilon_r>\varepsilon^\star\}$ and $\mathrm{prec}_{\mathrm{prop}}=1-|\mathcal{R}_{\mathrm{bad}}|/\hat M_{\mathrm{eff}}$. The expected number of members of $\mathcal{R}_{\mathrm{bad}}$ surviving $N_{\min}$ tests satisfies
$$\E[|\{r\in\mathcal{R}_{\mathrm{bad}}:\text{survives}\}|]\le (1-\mathrm{prec}_{\mathrm{prop}})\,\hat M_{\mathrm{eff}}\, e^{-\varepsilon^\star N_{\min}} \le \hat M_{\mathrm{eff}}\, e^{-\varepsilon^\star N_{\min}}.$$
\end{corollary}
\begin{proof}
Linearity of expectation over survival indicators for $r\in\mathcal{R}_{\mathrm{bad}}$, each bounded by $(1-\varepsilon_r)^{N_{\min}}\le(1-\varepsilon^\star)^{N_{\min}}\le e^{-\varepsilon^\star N_{\min}}$ since $\varepsilon_r>\varepsilon^\star$ on $\mathcal{R}_{\mathrm{bad}}$; $|\mathcal{R}_{\mathrm{bad}}|=(1-\mathrm{prec}_{\mathrm{prop}})\hat M_{\mathrm{eff}}$ by definition of $\mathrm{prec}_{\mathrm{prop}}$.
\end{proof}
These two readings are consistent, not opposed: the first bounds \emph{per-rule} certification cost and is cheap in $\hat M_{\mathrm{eff}}$; the second bounds the \emph{total count} of bad survivors at fixed cost and is expensive in $\hat M_{\mathrm{eff}}$ \emph{unless} the proposer's precision is high. Concretely, at $N_{\min}=200,\varepsilon^\star=0.01$: a proposer with $\mathrm{prec}_{\mathrm{prop}}=0$ (worst case) gives $\approx1.3\times10^5$ expected false survivors at $\hat M_{\mathrm{eff}}=10^6$ versus $13$ at $\hat M_{\mathrm{eff}}=10^2$; but a proposer with $\hat M_{\mathrm{eff}}=10^6$ and $\mathrm{prec}_{\mathrm{prop}}=0.9999$ controls this quantity just as well as the small, imprecise one. What justifies a parsimonious proposer is therefore the \emph{measured precision} of its proposals (Table 2 in the main paper), not the raw candidate count $\hat M_{\mathrm{eff}}$ taken alone --- and both quantities, $\hat M_{\mathrm{eff}}$ and $\mathrm{prec}_{\mathrm{prop}}$, are directly measurable in the P1/P2/P3 ablation (\autoref{app:protocol}).

\begin{theorem}[Chaining to rollout fidelity]
\label{thm:t3-full}
Under hard injection (projection or residual), disjoint masks, and assumption (H) that the rollout stays within the support of $d^\pi$: the one-step squared prediction error satisfies $\E\|\hat\omega_{t+1}-\omega_{t+1}\|^2 \le \varepsilon_d B^2+e_u$, and under an $L$-Lipschitz rollout operator, $E_H\le\big(\sum_{i=0}^{H-1}L^i\big)(\varepsilon_d B^2+e_u)$.
\end{theorem}
\begin{proof}
Orthogonal decomposition: $\|\hat\omega-\omega\|^2=\|P_d(\hat\omega-\omega)\|^2+\|P_u(\hat\omega-\omega)\|^2$. Under hard injection, the covered term is zero except on a violation event of probability $\le\varepsilon_d$, on which the error is bounded by $B^2$; hence $\E\|P_d(\cdot)\|^2\le\varepsilon_d B^2$. The residual term is $e_u$ by definition. Standard Lipschitz telescoping gives the $H$-step form.
\end{proof}
\textbf{Assumption (H) is false in general.} In open-loop rollout the model visits states it generated itself, outside $d^\pi$'s support. One can argue that hard projection keeps the rollout near the manifold of valid states, hence near the support --- but this argument is partially circular, and we state it as a conjecture, not a theorem. We recommend measuring rollout-support drift empirically (e.g., estimated density of visited states relative to $d^\pi$) rather than assuming (H).

\begin{theorem}[Degradation under policy shift]
\label{thm:t4-full}
Let $C_{\pi'/\pi}(c)=\sup_{x^-\in\{c\}}d^{\pi'}(x^-)/d^\pi(x^-)$. Then $\Pr_{\pi'}[V_r=1\wedge c]\le C_{\pi'/\pi}(c)\cdot\Pr_\pi[V_r=1\wedge c]$, and equivalently $\varepsilon^{\pi'}_r\le C_{\pi'/\pi}(c)\cdot\varepsilon^\pi_r\cdot\Pr_\pi[c]/\Pr_{\pi'}[c]$.
\end{theorem}
\begin{proof}
$\Pr_{\pi'}[V\wedge c]=\sum_{x^-\in\{c\}}d^{\pi'}(x^-)\Pr[V\mid x^-]\le C\sum_{x^-\in\{c\}}d^\pi(x^-)\Pr[V\mid x^-]=C\Pr_\pi[V\wedge c]$. The conditional form follows by dividing by $\Pr_{\pi'}[c]$ and rewriting $\Pr_\pi[V\wedge c]=\varepsilon^\pi_r\Pr_\pi[c]$.
\end{proof}
We recommend reporting the absolute form in the main text (``the rate of harmful triggers grows by at most a factor $C$''), which is directly interpretable, rather than the conditional form, whose ratio $\Pr_\pi[c]/\Pr_{\pi'}[c]$ is harder to read.

\begin{theorem}[Non-identifiability under partial observability]
\label{thm:t5-full}
There exist two Dec-POMDPs $\mathcal{M},\mathcal{M}'$, a policy $\pi$, an alternative policy $\pi'$, and a rule $r$ such that $\varepsilon^\pi_r=0$ in both $\mathcal{M}$ and $\mathcal{M}'$, yet $\varepsilon^{\pi'}_r>0$ in $\mathcal{M}$ and $\varepsilon^{\pi'}_r=0$ in $\mathcal{M}'$. Consequently, no certification procedure with access only to $d^\pi$ can distinguish $\mathcal{M}$ from $\mathcal{M}'$, nor bound $\varepsilon^{\pi'}_r$, for any $N_c$.
\end{theorem}
\begin{proof}[Construction and proof]
\textbf{Note on this construction (resolved per \texttt{revision/QUESTIONS.md}, item A5.4).} An earlier version of this proof used ``agent 2's position'' as the confound while agent 2 was implicitly a second decision-making agent of the Dec-POMDP. Since a partial invariant's reading $\varphi$ ranges over the \emph{joint} observation $\omega^j$ (main paper, Definition of a partial invariant), and since the rest of this paper's running examples treat an agent's own position as a standard component of its own observation, that version left it ambiguous --- and, read literally, more likely false --- whether agent 2's position was truly outside $\omega^j$, since agent 2 observing its own position would place it inside $\omega^2\subset\omega^j$ and hence inside the reading procedure's reach. We remove the ambiguity by making the confound an entity with \emph{no} observation channel at all, rather than by asserting a restriction on which agent sees what.

Let the state space be a four-cell corridor $\{0,1,2,3\}$ for a single decision-making agent, together with an exogenous hazard whose position $y_t\in\{1,2\}$ is part of the global state $s_t$ but is emitted by no agent's observation function: $y_t\notin\Omega^i$ for every $i$, and $y_t$ evolves according to a fixed environment process that is not itself a Dec-POMDP agent (it has no observation, no policy, and no reward) --- e.g., a moving obstacle on a fixed or randomized schedule, exactly as GridCraft's own moving hazards or other agents' \emph{scripted} (non-observing, non-learning) components would be modeled. The agent's observation is its own position $x^1_t$ only. Actions: \textsc{Right}, \textsc{Stay}.

Model $\mathcal{M}$: \textsc{Right} moves the agent one cell, \emph{unless} the target cell equals $y_t$, in which case the agent stays. Model $\mathcal{M}'$: \textsc{Right} always moves the agent one cell ($y_t$ has no causal effect on the transition).

Data-collection regime $\pi$: the hazard's exogenous process keeps $y_t=3$ throughout, while the agent's own policy visits $\{0,1,2\}$. Alternative regime $\pi'$: the hazard's exogenous process sometimes places $y_t\in\{1,2\}$ (e.g., a different level, schedule, or difficulty setting --- precisely the kind of shift gate F3, cross-variant falsification, is built to cover).

Candidate rule $r$: $a^1=\textsc{Right}\wedge x^1=k \Rightarrow x^1_{t+1}=k+1$, for $k\in\{0,1\}$.

\emph{Claim (i): the observation trajectories of the agent under $\pi$ are identically distributed in $\mathcal{M}$ and $\mathcal{M}'$.} Under $\pi$, $y_t$ never equals a cell in $\{1,2\}$, so the blocking mechanism of $\mathcal{M}$ is never triggered; the two transition kernels coincide on the support of $d^\pi$, and the agent's observations are therefore identically distributed under both models. Since $y_t\notin\omega^j$ for any agent, this identity holds regardless of what any reading $\varphi$ of the joint observation can access --- the confound is unavailable not because of a scoping choice on $\varphi$, but because it is not part of $\omega^j$ at all.

\emph{Claim (ii): $\varepsilon^\pi_r=0$ in both models.} Immediate from (i), since $r$ never sees a violation under $\pi$ in either model.

\emph{Claim (iii): $\varepsilon^{\pi'}_r>0$ in $\mathcal{M}$ and $\varepsilon^{\pi'}_r=0$ in $\mathcal{M}'$.} Under $\pi'$ in $\mathcal{M}$, the event ``$y_t=k+1$'' has non-zero probability and, by the blocking mechanism, produces a violation of $r$; in $\mathcal{M}'$ there is no blocking mechanism, so $\varepsilon^{\pi'}_r=0$ regardless of $y_t$.

The impossibility claim follows from (i): two models that are indistinguishable in the distribution they induce over any measurable function of $d^\pi$ (in particular, over any certification procedure's output, since that output can only be a function of $\omega^j$ and hence never of $y_t$) must yield the same output for both models; since the correct answer to ``is $r$ valid under $\pi'$'' differs between $\mathcal{M}$ and $\mathcal{M}'$, no procedure using only $d^\pi$ can be correct for both.
\end{proof}

\begin{corollary}[Memory is not a remedy]
Increasing the memory order $m$ in Family B transports (\autoref{app:dsl}) does not resolve this case: $y_t$ never appears, at any lag, in any agent's observation history, by construction ($y_t\notin\Omega^i$ for all $i$), not merely by omission.
\end{corollary}

This construction separates two failure modes that intuition often conflates:
\begin{center}
\begin{tabular}{@{}p{3.1cm}p{3.6cm}p{4.2cm}@{}}
\toprule
Nature & Where the information is & Remedy \\
\midrule
Deferred effect & In $h_{t-m}$, for some finite $m$ & Increase $m$ (memory-order test) \\
Latent confounder & Never in the observation history & Falsify against multiple policies/variants (F2--F4) \\
\bottomrule
\end{tabular}
\end{center}

\begin{theorem}[Sufficient condition for transfer]
\label{thm:t7-full}
Let $r$ be certified against each of $\pi_1,\dots,\pi_K$ (the falsification-gate policies of \autoref{thm:t2-full}), i.e., $\varepsilon^{\pi_k}_r\le\varepsilon$ for every $k$, and let $\bar\pi$ be their mixture: $d^{\bar\pi}=\frac1K\sum_k d^{\pi_k}$. If a deployment policy $\pi'$ satisfies $C_{\pi'/\bar\pi}(c):=\sup_{x^-\in\{c\}}d^{\pi'}(x^-)/d^{\bar\pi}(x^-)<\infty$, then
$$\varepsilon^{\pi'}_r \le C_{\pi'/\bar\pi}(c)\cdot\max_k\varepsilon^{\pi_k}_r.$$
\end{theorem}
\begin{proof}
First bound the mixture-level rate: $\Pr_{\bar\pi}[V_r=1\wedge c] = \frac1K\sum_k\Pr_{\pi_k}[V_r=1\wedge c] = \frac1K\sum_k \varepsilon^{\pi_k}_r\Pr_{\pi_k}[c] \le \big(\max_k\varepsilon^{\pi_k}_r\big)\cdot\frac1K\sum_k\Pr_{\pi_k}[c] = \big(\max_k\varepsilon^{\pi_k}_r\big)\Pr_{\bar\pi}[c]$, i.e., $\varepsilon^{\bar\pi}_r\le\max_k\varepsilon^{\pi_k}_r$ (a rule certified against every member of a family is certified, at least as well, against any mixture of that family --- a convexity argument, not an additional assumption). Now apply \autoref{thm:t4-full} with $\bar\pi$ in place of $\pi$: $\Pr_{\pi'}[V_r=1\wedge c]\le C_{\pi'/\bar\pi}(c)\Pr_{\bar\pi}[V_r=1\wedge c]\le C_{\pi'/\bar\pi}(c)\big(\max_k\varepsilon^{\pi_k}_r\big)\Pr_{\bar\pi}[c]$, and dividing by $\Pr_{\pi'}[c]$ gives the conditional form stated (up to the same $\Pr_{\bar\pi}[c]/\Pr_{\pi'}[c]$ ratio discussed after \autoref{thm:t4-full}; we again recommend reporting the absolute form $\Pr_{\pi'}[V_r=1\wedge c]\le C_{\pi'/\bar\pi}(c)\max_k\varepsilon^{\pi_k}_r\Pr_{\bar\pi}[c]$ in practice).
\end{proof}
This is precisely what gates F2, F3, and F4 are designed to produce: falsifying against an active-exploration policy, environment variants, and alternative generators is exactly falsifying against a family $\{\pi_k\}$ rather than a single $\pi$, and \autoref{thm:t7-full} says this family's \emph{coverage} of the deployment policy's precondition region (a finite $C_{\pi'/\bar\pi}(c)$) is what licenses transfer. \autoref{thm:t5-full} shows this coverage condition is not automatic --- a single passive policy can have $C_{\pi'/\pi}(c)=\infty$ for a $\pi'$ that only \autoref{thm:t7-full}'s broader family would cover --- so the negative result does not invalidate the method; it identifies exactly which part of the protocol (active, multi-policy, multi-variant falsification) is load-bearing, and \autoref{thm:t7-full} is the corresponding positive guarantee once that part is in place.

\begin{remark}[Composition, full derivation]
If rules $r_1,\dots,r_p$ are injected on disjoint feature positions, the aggregate covered violation rate satisfies $\varepsilon_d\le\sum_{i=1}^p\varepsilon_{r_i}$ by a union bound over the individual violation events. This means coverage is not free: injecting $50$ rules at $\varepsilon=0.01$ each already gives $\varepsilon_d\le0.5$. The hard-injection threshold $\tau_{\mathrm{hard}}$ should therefore be calibrated against a total error budget (e.g., $\tau_{\mathrm{hard}}$ adjusted to a target $1-\varepsilon_{\mathrm{cible}}/p$), not fixed per rule independent of how many rules are simultaneously active.
\end{remark}

\section{Extended experimental protocol}
\label{app:protocol}

\subsection{Environments and difficulty gradient}
GridCraft, Overcooked, Predator--Prey, and SMACv2, as in prior NS-MAWM work. On GridCraft, invariants are organized along a difficulty gradient designed to test \autoref{thm:t5-full} empirically: spatial/boundary rules $\to$ persistence $\to$ conservation $\to$ long-memory cooldowns $\to$ rules over a planted, latent environment hazard never entering any agent's observation (predicted uncertifiable under F1-only falsification, and predicted to become certifiable once F2--F4 provide the coverage \autoref{thm:t7-full} requires).

\subsection{Anonymization protocol}
Permutation of feature indices; stripping of semantic names; injection of distractor features (irrelevant noise) and constant features (to test the non-vacuity screen); rescaling. We deliberately do \emph{not} apply an arbitrary linear mixing of coordinates, since that would make invariants inexpressible in the catalog of \autoref{app:dsl}; this is a real dependency on a reasonably factorized representation, reported as a limitation rather than hidden.

\subsection{Metrics}
\begin{itemize}
  \item Fraction recovered $=(M-B_1)/(B_2-B_1)$, where $B_1$ = purely neural MAWM, $B_2$ = NS-MAWM with hand-written oracle rules, $M$ = NS-MAWM with discovered rules.
  \item Precision/recall of discovered rules against the oracle rule set (already written for GridCraft in prior work, used here as ground truth for the discovery metric only, not injected into $M$).
  \item $\hat\kappa$ vs.\ oracle $\kappa$; false-discovery rate; retrogradation rate; discovery wall-clock time.
  \item $\mathrm{RVR}_{\mathrm{pre}}$, $\mathrm{RVR}_{\mathrm{post}}$, correction magnitude $\|P_d\hat\omega^{pre,j}_{t+1}-\omega^{d,j}_{t+1}\|$.
  \item Residual-prediction degradation when covered features correlated with the residual target are deliberately masked (the correlation ablation).
  \item Certification outcome per difficulty-gradient class, reported separately for F1-only vs.\ F1--F4 falsification.
\end{itemize}
All quantities reported as mean $\pm$ standard error across at least 5 independent seeds; paired comparisons for close differences; overlapping confidence intervals are not treated as a strict ranking.

\subsection{Ablations}
Proposer (P1 vs.\ P2 vs.\ P3); falsification severity (F1 vs.\ F1+F2 vs.\ F1--F4), measuring false-discovery rate at each level; memory order $m_{\max}=0$ vs.\ $>0$; non-vacuity screen on vs.\ off; soft-only vs.\ hard-only vs.\ graded promotion; retrogradation on vs.\ off; precondition depth $\beta=2$ vs.\ $3$; residual-head conditioning on vs.\ off $\omega^{d,j}_{t+1}$ (isolating the fix in the main paper's injection section).

\subsection{Result tables (placeholders)}
The main paper's discovery-vs-oracle and difficulty-gradient tables specify the exact column structure to fill from seed-level logs. No additional table is introduced here; this document specifies protocol, not results.

\section{LLM-proposer prompting protocol}
\label{app:llm}

The P3 proposer receives: (i) the anonymized feature schema (index, type, observed value range, no semantic name); (ii) the action space; (iii) a sample of example transitions drawn from $\mathcal{D}_{\mathrm{prop}}$; and is asked to return candidates strictly within the DSL of \autoref{app:dsl} (a reading, a precondition of depth $\le\beta$, and a transport from Families A/B/C), in a structured, machine-parseable schema. Candidates that do not parse into the DSL are discarded before falsification and do not count toward $\hat M$; candidates that parse but are semantically invalid (e.g., a transport referencing an undefined reading) are discarded at the same stage. We report the syntactic-validity rate as a diagnostic of the proposer, separate from the falsification results.

\section{Reproducibility notes}
\label{app:repro}

Discovery introduces additional artifacts beyond the injection-time reproducibility material already used in prior NS-MAWM work: the four data partitions ($\mathcal{D}_{\mathrm{prop}},\mathcal{D}_{\mathrm{fit}},\mathcal{D}_{\mathrm{test}},\mathcal{D}_{\mathrm{shift}}$) and their exact construction per environment; the enumerated catalog and its realized $\hat M$ for P1; the decision-tree hyperparameters for P2; the exact LLM prompts, model identifier, and sampling temperature for P3; and the discovered rule inventory $\{(\varphi,c,\rho,m,q,\mathrm{cov},N_c)\}$ per seed, which should be logged in full (not only aggregated) so that precision/recall against the oracle set can be recomputed independently.

\bibliographystyle{plainnat}
\bibliography{references}

\end{document}
